\chapter{Conclusion}
\label{chapter6}

This thesis presented a hardware--software co-design for securing embedded programming and
execution under physical-access threat models. The approach introduces a minimal ISA extension
and a custom authorization instruction (\texttt{kc}) that triggers verification against key
material stored in protected non-volatile memory. A constant-time comparator and an
Authenticator FSM gate both execution and debug enablement, ensuring that privileged capabilities
are released only after successful, device-local authorization. Integration into the toolchain
allows the authorization primitive to be emitted and preserved through compilation, supporting
reproducible and auditable binaries.

The implementation demonstrated end-to-end enforceability across the host workflow and device
logic. Toolchain artifacts confirmed that the custom instruction is preserved through the build
pipeline, and host-side logs showed that key exchange precedes debug enablement. Programming
workflow evidence demonstrated connection and image load under the intended authorization policy.
Collectively, these results support the practicality of instruction-triggered, key-gated control
for restricting unauthorized reflashing and invasive debugging in embedded systems.

Despite the encouraging results, several limitations remain. The evaluation emphasized functional
correctness and workflow evidence rather than a comprehensive characterization of hardware
overheads and latency. The key verification datapath currently uses a single comparison strategy,
and the impact of fault/glitch or side-channel attacks has not been quantified. In addition, the
authorization protocol was evaluated on a representative workflow; broader deployment requires
formalized provisioning procedures, recovery mechanisms, and lifecycle policy integration.

Future work should therefore quantify area, power, and timing overheads across target devices,
extend the protocol to include key rotation and revocation, and harden the verification path
against fault injection and side-channel leakage. Additional toolchain regression tests and
formal modeling of the FSM would further raise assurance. Finally, multi-device deployment
studies with centralized audit policy would validate scalability and operational robustness.

In summary, this thesis establishes that device-local authorization can be enforced via a minimal
ISA extension and protected key storage, providing a practical foundation for securing embedded
programming and debug access.
